\documentclass[11pt]{article}
\usepackage{amsmath,amssymb,amsthm}
\usepackage[margin=1in]{geometry}
\newtheorem{lemma}{Lemma}
\newtheorem{theorem}[lemma]{Theorem}
\newtheorem{corollary}[lemma]{Corollary}
\newcommand{\Stab}{\operatorname{Stab}}
\newcommand{\rad}{\operatorname{rad}}
\title{Involutions preserving an alternating form in dimension seven}
\author{Research notes for the F004 proof programme}
\date{}
\begin{document}
\maketitle

Let $F$ be a field of characteristic different from $2$, let $V$ have
dimension $7$ over $F$, and let $\varphi\in\Lambda^3V^*$. Its contraction
form is the symmetric bilinear map
\[
 b_\varphi:V\times V\longrightarrow\Lambda^7V^*,\qquad
 b_\varphi(x,y)=(\iota_x\varphi)\wedge(\iota_y\varphi)\wedge\varphi.
\]
Throughout this note it is nondegenerate, meaning that its radical is zero.
Write $H=\Stab_{\operatorname{GL}(V)}(\varphi)$. No identification of $H$
with a named algebraic or finite group is assumed.

\begin{lemma}\label{lem:covariance}
For $h\in H$ and $x,y\in V$,
\begin{equation}\label{eq:covariance}
 b_\varphi(hx,hy)=\det(h)^{-1}b_\varphi(x,y).
\end{equation}
Consequently $\det(h)^9=1$. In particular, an involution in $H$ has
determinant $1$.
\end{lemma}

\begin{proof}
The equality $h^*\varphi=\varphi$ implies
\[
 \iota_{hx}\varphi=(h^{-1})^*(\iota_x\varphi).
\]
Indeed, evaluating the right side at $u,v\in V$ gives
$\varphi(x,h^{-1}u,h^{-1}v)=\varphi(hx,u,v)$. Since $h^{-1}$ also
preserves $\varphi$, naturality of the exterior product now gives
\[
 b_\varphi(hx,hy)=(h^{-1})^*b_\varphi(x,y).
\]
The action on $\Lambda^7V^*$ is multiplication by $\det(h^{-1})$, proving
\eqref{eq:covariance}.

Choose a nonzero volume form $\Omega$ and write
$b_\varphi(x,y)=B(x,y)\Omega$. In any basis the matrix $B$ is nonsingular.
If $A$ is the matrix of $h$, equation~\eqref{eq:covariance} becomes
\[
 A^{\mathsf T}BA=\det(A)^{-1}B.
\]
Taking determinants and cancelling $\det(B)\ne0$ gives
$\det(A)^2=\det(A)^{-7}$, hence $\det(A)^9=1$. If $h^2=I$, its
determinant also has square $1$. Therefore
\[
 \det(h)=\frac{\det(h)^9}{(\det(h)^2)^4}=1.
\]
The conclusion does not require choosing a preferred volume form.
\end{proof}

\begin{lemma}\label{lem:negative-two}
An involution in $H$ cannot have a two-dimensional $-1$ eigenspace.
\end{lemma}

\begin{proof}
Suppose otherwise. Let $V=L\oplus W$ be its $+1$ and $-1$ eigenspaces, of
dimensions $5$ and $2$. The direct sum follows from the characteristic
assumption. Choose a basis $\alpha,\beta$ of $W^*$ and extend it to vanish
on $L$. Forms on $L$ are similarly extended to vanish on $W$.

Invariance of $\varphi$ removes the summands containing an odd number of
covectors on $W$. Since $\dim W=2$, the remaining terms have the unique
shape
\[
 \varphi=\eta+\gamma\wedge\alpha\wedge\beta,
 \qquad \eta\in\Lambda^3L^*,\quad\gamma\in L^*.
\]
For $w\in W$, contraction therefore gives
\[
 \iota_w\varphi
 =-\gamma\wedge\bigl(\alpha(w)\beta-\beta(w)\alpha\bigr).
\]
Thus, for any $w,z\in W$, the product
$(\iota_w\varphi)\wedge(\iota_z\varphi)$ contains $\gamma$ twice and
vanishes. It follows directly that $b_\varphi(W,W)=0$.

By Lemma~\ref{lem:covariance} the involution has determinant $1$, so it
preserves $b_\varphi$. For $w\in W$ and $\ell\in L$ this yields
\[
 b_\varphi(w,\ell)=b_\varphi(-w,\ell)=-b_\varphi(w,\ell).
\]
Since $2\ne0$ in $F$, we obtain $b_\varphi(W,L)=0$ as well. Consequently
$W\subseteq\rad b_\varphi$, contradicting nondegeneracy.
\end{proof}

\paragraph{Imported radical lemma.}
We use Theorem~1 of the companion note \emph{A radical forced by the
eigenspaces of an involution},
\texttt{alternating\_form\_involution\_obstruction.tex}, exact edition
7060 bytes with SHA256
\texttt{247875d240fe8529c6ae5505febaf26d4dd9199037869325af548e49c0ddd54f}.
It says that a preserving involution with eigenspace dimensions $1$ and $6$
for $+1$ and $-1$ has its entire negative eigenspace in $\rad b_\varphi$.
Its explicit contraction calculation is reused here, rather than repeated.
This is an internal proof dependency; its statement and proof must be
reviewed with this note.

\begin{theorem}\label{thm:fixed-three}
Every nonidentity involution in $H$ has a three-dimensional $+1$ eigenspace
and a four-dimensional $-1$ eigenspace. Its trace on $V$ is $-1$.
\end{theorem}

\begin{proof}
Let $s\in H$ satisfy $s^2=I$ and $s\ne I$. The polynomial
$(t-1)(t+1)$ has distinct roots in $F$, so $V$ is the direct sum of the two
eigenspaces. Write $q$ for the dimension of the negative eigenspace.
Lemma~\ref{lem:covariance} gives $(-1)^q=\det(s)=1$. Since the
characteristic is not $2$, $q$ is even, and thus
\[
 q\in\{0,2,4,6\}.
\]
The value $q=0$ would give $s=I$, contrary to the hypothesis. The value
$q=2$ is excluded by Lemma~\ref{lem:negative-two}. The value $q=6$ is
excluded by the imported radical lemma and nondegeneracy of $b_\varphi$.
Hence $q=4$ and the positive dimension is $7-4=3$. The trace is therefore
$3-4=-1$, interpreted as an element of $F$.
\end{proof}

\begin{corollary}\label{cor:idempotent}
If $E\in\operatorname{End}_F(V)$ satisfies $E^2=E$ and
$(2E-I)^*\varphi=\varphi$, then $\operatorname{rank}E\in\{3,7\}$. In
the rank $7$ case $E=I$ and $2E-I=I$.
\end{corollary}

\begin{proof}
The idempotent gives $V=\operatorname{im}E\oplus\ker E$. The map $2E-I$
is an involution with positive eigenspace $\operatorname{im}E$ and negative
eigenspace $\ker E$. If it is not the identity,
Theorem~\ref{thm:fixed-three} gives $\dim\operatorname{im}E=3$. Otherwise
$2E=2I$, and the characteristic assumption gives $E=I$.
\end{proof}

\paragraph{Reusable exclusion.}
If $G\leq H$, an actual nonidentity involution with fixed-space dimension
different from $3$ is not in $G$. Moreover, if $D\leq G$, neither coset
$Ds$ nor $sD$ of such an excluded element can meet $H$: otherwise multiplying
by the inverse of the corresponding element of $D$ would put $s$ in $H$.
This reuses the subgroup argument in Proposition~4 of the companion note.
No assumption that $s$ normalizes or centralizes $D$ is required.

\paragraph{Attribution and application scope.}
The multiplicity argument arose in the Root mathematical discussion. Its
ingredients are standard exterior algebra, determinant identities and the
elementary eigenspace decomposition of an involution. No claim of historical
novelty, group conjugacy classification, or priority is made. The theorem
identifies necessary eigenspace multiplicities; it does not say that all
involutions of that type are conjugate, or that the type is sufficient for
preserving a particular form.

Applying the theorem to a named matrix group requires the exact imported
identification of its seven-dimensional model, the actual preserved form and
the nondegenerate contraction form. Applying its projector corollary also
requires the actual idempotent and the exact identity $s=2E-I$. Those
identifications have not been established for position2739 by this note.
The existing numerical end-intersection certificate and accepted case proof
are unchanged. These notes supply a possible shared proof leg, not another
scientific execution, Original row, conjecture answer or manuscript recovery.

\end{document}
